\section{Hazard identification}
\label{chapter4}
 
The only source of physical harm in this system is the rotating actor (rotor).
The safety function of the demonstrator is: if an object is closer to the rotor
than the speed-dependent minimum distance, the motor must be switched off.
The hazards below are therefore all situations in which this safety function is
lost, is applied too late, or is based on wrong data.
 
\subsection{Identified hazards and countermeasures}
 
	\begin{enumerate}
		\item \textbf{Hazard 1: Undetected failure of a distance sensor} \\
			\emph{Cause:} An HC-SR04 delivers no echo at all (broken wiring, defective
			transducer), returns a permanently frozen value, returns a value outside the
			specified measuring range, or receives the echo of the second sensor
			(cross-talk). \\
			\emph{Effect:} The control unit keeps working with an invalid or outdated
			distance value and interprets the danger zone as free. The motor is not
			switched off although a person has already entered the protected area.
			\\
			\emph{Countermeasure in software (Raspberry Pi 2):}
			\begin{itemize}
				\item Echo timeout monitoring: if no echo edge is received within the
					maximum expected time, the measurement is marked as invalid. An
					absent echo is never interpreted as ``no object''.
				\item Range check: measurements outside the specified measuring range
					of the sensor are rejected.
				\item Stuck-value detection: if a sensor reports an unchanged value over
					a defined number of consecutive cycles, it is treated as faulty.
				\item Plausibility check between both sensors: if the difference of the
					two measurements exceeds a defined limit, no trustworthy distance
					information exists.
				\item Age check: every distance value carries a timestamp. A value older
					than one monitoring cycle is discarded and never reused. No distance
					variable is evaluated before a valid measurement has been written to it.
				\item Time multiplexing: only one sensor is triggered at a time, so an
					echo can always be assigned to exactly one sensor.
				\item Fail-safe reaction: each of the checks above leads to the safe
					state (motor off). The error state is latched and displayed; the
					system restarts only after a manual acknowledgement.
			\end{itemize}
			Requirement see \ref{req.2}
		\item \textbf{Hazard 2: Protection zone too small for the actual rotor speed} \\
			\emph{Cause:} The minimum distance is derived from the speed setpoint that
			is entered on the host unit and transmitted via Bluetooth. The actual rotor
			speed can deviate from this value during run-up, under changing load, after
			a corrupted or outdated transmission, or because the rotor continues to
			coast down after the motor has been switched off. \\
			\emph{Effect:} The monitored area is smaller than required for the real
			speed, or the system reports the safe state while the rotor is still
			rotating. A person can reach the rotor while it still carries dangerous
			energy. \\
			\emph{Countermeasure in software (Raspberry Pi 2):}
			\begin{itemize}
				\item The minimum distance is calculated from the measured rotor speed,
					not from the transmitted setpoint. If no valid speed measurement is
					available, the maximum speed is assumed as the worst case.
				\item The minimum distance includes the reaction time of the monitoring
					cycle and the coast-down time until standstill, not only the
					geometric distance to the rotor.
				\item Setpoints received via Bluetooth are accepted only with a valid
					checksum, within the allowed value range and with a current
					timestamp. An outdated setpoint is discarded.
				\item A speed increase is only executed after the correspondingly larger
					protection zone has been verified as free.
				\item The safe state is signalled only when the measured speed is zero,
					not when the switch-off command has been issued.
			\end{itemize}
			Requirement see \ref{req.3}
		\item \textbf{Hazard 3: Loss of cyclic monitoring while the motor is running} \\
			\emph{Cause:} Software fault on the control unit: endless or blocked loop
			(for example waiting for an echo that never arrives), deadlock, crash of the
			control process, or an exceeded cycle time caused by system load. \\
			\emph{Effect:} The motor keeps its last enable state while the protected
			area is no longer monitored. The safety function is lost silently, without
			any visible reaction of the system. \\
			\emph{Countermeasure in software (Raspberry Pi 2):}
			\begin{itemize}
				\item Hardware watchdog that is retriggered only at the end of a
					completely executed monitoring cycle. An expired watchdog resets the
					controller; the enable output then defaults to ``off''.
				\item Dynamic enable signal: the motor driver is enabled by a periodic
					signal that only a running and complete cycle can generate. A static
					signal level cannot keep the motor switched on.
				\item Cycle time monitoring against a defined deadline. A missed
					deadline leads to the safe state.
				\item All read operations on sensors and on the communication interface
					use timeouts. There are no unbounded blocking waits in the safety
					related path.
				\item Self-test of the complete safety path at program start (sensor
					reading, enable output, error indication) before the motor can be
					started for the first time.
			\end{itemize}
			Requirement see \ref{req.1}
	\end{enumerate}
 
\subsection{Identified hazards without countermeasures }
 
	\begin{enumerate}
		\item \textbf{Hazard 1: Mechanical failure of the rotor} \\
			\emph{Cause:} Material fatigue, a loose fastening, imbalance at high speed
			or a foreign object entering the rotor. \\
			\emph{Effect:} Parts are ejected with high energy and can cause injuries
			also outside the monitored area. \\
			\emph{Reason why no countermeasure is implemented:} The event occurs within
			milliseconds and cannot be detected by distance sensors. Software can
			neither prevent it nor reduce its consequences. Only mechanical measures
			such as an enclosure, a defined mounting torque and a limited maximum speed
			are effective, and these are outside the software scope of this project. \\
			\emph{Residual risk:} Accepted and documented. The demonstrator is operated
			with a limited maximum speed.
		\item \textbf{Hazard 2: Failure of the motor driver (switch-off command without effect)} \\
			\emph{Cause:} A welded or short-circuited switching element in the power
			stage, a defective enable input, or a driver that remains supplied although
			the control unit commands the switch-off. \\
			\emph{Effect:} The motor keeps running although the software has correctly
			detected the hazard and has commanded the shutdown. A single hardware fault
			bypasses the complete safety function. \\
			\emph{Reason why no countermeasure is implemented:} The software cannot
			detect or override a switching element that is stuck in the conducting
			state. A second, independent shutdown path would be required, for example a
			safety relay in the supply line with feedback contact monitoring, or a
			driver with a dedicated safe torque off input. The demonstrator uses a
			single-channel shutdown path; adding redundant hardware is outside the
			scope of this project. \\
			\emph{Residual risk:} Documented single point of failure.
		\item \textbf{Hazard 3: Physical detection limits of the ultrasonic sensors} \\
			\emph{Cause:} The HC-SR04 evaluates reflected sound and covers only a narrow
			detection cone. Objects approaching from the side, from above or from below
			are outside this cone. Sound absorbing materials such as soft clothing,
			gloves or foam, and strongly inclined surfaces that deflect the echo,
			produce no usable reflection. Small objects such as fingers, an approach
			faster than the measurement cycle, other ultrasonic sources and strong air
			movement further reduce the detection reliability. \\
			\emph{Effect:} A person can enter the danger zone without being detected.
			The motor is not switched off, although the software works correctly. \\
			\emph{Reason why no countermeasure is implemented:} These are properties of
			the measuring principle and not software faults. No plausibility check can
			detect an object that produces no echo at all. A reliable solution would
			require a certified safety sensor with a defined detection capability and
			internal self-monitoring, for example a safety laser scanner or a light
			curtain. The HC-SR04 is a standard component without any safety rating. \\
			\emph{Residual risk:} The demonstrator shows the functional principle only.
			It is explicitly not suitable for the protection of persons.
	\end{enumerate}

